% !TEX program = lualatex
% Práctica guiada de Fundamentos de las Matemáticas I
% Lenguaje matemático y definiciones: categoría de lenguajes naturales

\documentclass[11pt,a4paper]{article}

\usepackage{fontspec}
\usepackage[spanish,es-nodecimaldot]{babel}
\usepackage[a4paper,margin=2.7cm]{geometry}
\usepackage{amsmath,amssymb,mathtools}
\usepackage{enumitem}
\usepackage{xcolor}
\usepackage{booktabs}
\usepackage{hyperref}

\setmainfont{Latin Modern Roman}
\setsansfont{Latin Modern Sans}
\setmonofont{Latin Modern Mono}

\definecolor{FdMBlue}{RGB}{20,55,100}
\definecolor{FdMGray}{RGB}{245,245,245}
\definecolor{FdMRed}{RGB}{130,30,30}
\definecolor{FdMGreen}{RGB}{25,100,60}

\hypersetup{
  colorlinks=true,
  linkcolor=FdMBlue,
  urlcolor=FdMBlue
}

\newcommand{\Cat}{\mathcal{L}}
\newcommand{\Ob}{\operatorname{Ob}}
\newcommand{\Hom}{\operatorname{Hom}}
\newcommand{\id}{\operatorname{id}}

\setlist[itemize]{leftmargin=2em}
\setlist[enumerate]{leftmargin=2.5em}

\title{\textbf{Práctica guiada 1}\\
Construcción idealizada de una categoría de lenguajes naturales}
\author{Fundamentos de las Matemáticas I}
\date{}

\begin{document}

\maketitle

\section*{Objetivo}

El objetivo de esta práctica es utilizar una definición matemática no trivial
para entrenar las ideas del capítulo 1: declarar objetos, introducir símbolos,
fijar operaciones y comprobar reglas.

Trabajaremos con una construcción idealizada: una categoría cuyos objetos son
lenguajes naturales y cuyos morfismos representan traducciones.

\section*{Recordatorio: definición de categoría}

Una categoría \(\mathcal{C}\) consta de los siguientes datos:

\begin{itemize}
  \item una colección de objetos, denotada \(\Ob(\mathcal{C})\);
  \item para cada par de objetos \(A,B\), una colección de morfismos
        \(\Hom_{\mathcal{C}}(A,B)\);
  \item una operación de composición de morfismos;
  \item para cada objeto \(A\), un morfismo identidad \(\id_A\).
\end{itemize}

Estos datos deben cumplir dos reglas:

\begin{enumerate}
  \item \textbf{Asociatividad.} Si las composiciones tienen sentido, entonces
  \[
    h\circ (g\circ f)=(h\circ g)\circ f.
  \]

  \item \textbf{Identidades.} Si \(f:A\to B\), entonces
  \[
    f\circ \id_A=f
    \quad\text{y}\quad
    \id_B\circ f=f.
  \]
\end{enumerate}

\section*{Situación de partida}

Queremos construir una categoría \(\Cat\) de lenguajes naturales. Para que el
ejercicio sea manejable, trabajaremos con el siguiente conjunto de lenguajes:
\[
  \{\text{español},\text{inglés},\text{francés}\}.
\]

La idea intuitiva es que una traducción del español al inglés sea un morfismo
del objeto ``español'' al objeto ``inglés''.

\section*{Parte 1: objetos}

\begin{enumerate}
  \item Define explícitamente la colección de objetos de \(\Cat\):
  \[
    \Ob(\Cat)=\{\ldots\}.
  \]

  \item Explica en una frase qué representa cada objeto.

  \item Añade un cuarto lenguaje a la colección y escribe de nuevo
  \(\Ob(\Cat)\).
\end{enumerate}

\section*{Parte 2: morfismos}

Para cada par de lenguajes \(A\) y \(B\), escribiremos
\[
  \Hom_{\Cat}(A,B)
\]
para referirnos a las traducciones idealizadas de \(A\) a \(B\).

\begin{enumerate}
  \item Escribe qué debería representar el conjunto
  \[
    \Hom_{\Cat}(\text{español},\text{inglés}).
  \]

  \item Escribe qué debería representar el conjunto
  \[
    \Hom_{\Cat}(\text{francés},\text{español}).
  \]

  \item Si \(t\in \Hom_{\Cat}(\text{español},\text{inglés})\), completa la frase:

  \begin{quote}
    El símbolo \(t\) representa \ldots
  \end{quote}

  \item Explica por qué no debemos escribir simplemente \(t\) sin indicar antes
  en qué conjunto de morfismos vive.
\end{enumerate}

\section*{Parte 3: identidades}

Para cada lenguaje \(A\), necesitamos un morfismo identidad
\[
  \id_A:A\to A.
\]

\begin{enumerate}
  \item Describe con palabras qué debería ser
  \[
    \id_{\text{español}}.
  \]

  \item Describe con palabras qué debería ser
  \[
    \id_{\text{inglés}}.
  \]

  \item Explica por qué una identidad no debería cambiar el significado de un
  texto.
\end{enumerate}

\section*{Parte 4: composición}

Supongamos que tenemos dos traducciones idealizadas:
\[
  f:\text{español}\to \text{inglés},
  \qquad
  g:\text{inglés}\to \text{francés}.
\]

Entonces podemos formar la composición
\[
  g\circ f:\text{español}\to \text{francés}.
\]

\begin{enumerate}
  \item Explica qué significa traducir primero mediante \(f\) y después mediante
  \(g\).

  \item Completa:
  \[
    f\in \Hom_{\Cat}(\ldots,\ldots),
    \qquad
    g\in \Hom_{\Cat}(\ldots,\ldots),
  \]
  y
  \[
    g\circ f\in \Hom_{\Cat}(\ldots,\ldots).
  \]

  \item Decide si la composición siguiente tiene sentido:
  \[
    f\circ g.
  \]
  Justifica la respuesta indicando dominios y codominios.
\end{enumerate}

\section*{Parte 5: comprobación de axiomas}

Sean
\[
  f:\text{español}\to\text{inglés},
  \qquad
  g:\text{inglés}\to\text{francés},
  \qquad
  h:\text{francés}\to\text{alemán}.
\]

\begin{enumerate}
  \item Escribe los dominios y codominios de \(g\circ f\) y de \(h\circ g\).

  \item Explica con palabras qué significa
  \[
    h\circ(g\circ f).
  \]

  \item Explica con palabras qué significa
  \[
    (h\circ g)\circ f.
  \]

  \item En el modelo idealizado, ambas expresiones deberían representar la misma
  traducción compuesta. ¿Qué axioma de categoría se está comprobando?
\end{enumerate}

\section*{Parte 6: límites del modelo}

La construcción anterior es útil, pero las traducciones reales no se comportan
exactamente como funciones perfectas.

\begin{enumerate}
  \item Da un ejemplo de una palabra o expresión que pueda tener más de una
  traducción.

  \item Explica por qué una traducción puede depender del contexto.

  \item Explica qué puede perderse al traducir sucesivamente
  \[
    \text{español}\to\text{inglés}\to\text{francés}.
  \]

  \item Indica qué estamos idealizando cuando tratamos las traducciones como
  morfismos de una categoría.
\end{enumerate}

\section*{Entrega}

Entrega una respuesta breve, de una o dos páginas, que incluya:

\begin{enumerate}
  \item la definición de los objetos de \(\Cat\);
  \item la interpretación de los morfismos;
  \item la identidad de cada lenguaje;
  \item un ejemplo de composición de traducciones;
  \item una comprobación informal de asociatividad;
  \item una explicación de los límites del modelo.
\end{enumerate}

\section*{Pregunta final}

Explica en cinco o seis líneas por qué esta práctica ilustra el principio:

\begin{quote}
  Todo símbolo relevante debe definirse antes de utilizarse.
\end{quote}

\end{document}
